% Korean edition of sections/s2_3_queries.tex.
\seclabel
\claim{\secnum{2.3}{sec:queries}\quad 거래일이 아닌 날짜를 보내면 HTTP 200과 함께 다른 날의 데이터가 온다}

\lead{모든 REST 요청은 \code{/api/{area}/{name}} 형태이며, 경로는 로그인 뒤에만 내려오는 코드에서 얻었다.} 공개 JavaScript 번들에는 \code{Auth} 경로만 있다. 거래 화면의 코드는 \code{/assets/app/} 아래에 있고 \code{asset_token} 쿠키가 있어야 내려오므로(\figref{creds}), 나머지 경로는 로그인 뒤 내려받은 코드에서 얻었다(\figref{paths}). 관리자 영역(\code{Master}, \code{Admin})은 래퍼의 범위 밖이다.

% Korean edition of figs/paths.tex.
\fig{paths}{래퍼의 REST 경로는 5개 영역으로 나뉘며, 굵은 경로는 실제로 호출했다}[웹 클라이언트 코드에서 뽑은 주요 REST 경로, 경로의 첫 부분으로 묶음]{%
{\tabfont\begin{tabularx}{\linewidth}{>{\RaggedRight\arraybackslash}p{16mm} >{\RaggedRight\arraybackslash}X}
\thd{영역} & \thd{경로 (일부)}\\
\code{Auth} & \textbf{\code{Login}}, \textbf{\code{SilentRefresh}}, \textbf{\code{Logout}}\\\rowsep
\code{Portfolio} & 주문: \textbf{\code{AddOrder}}, \code{ModifyOrder}, \code{CancelOrder}, \code{DeleteOrder}, \textbf{\code{CancelOrderBatch}}, \textbf{\code{ValidateProduct}}, \code{OrderLogs}, \textbf{\code{Orders}}\\
 & 상태: \textbf{\code{PfBs}}, \textbf{\code{Session}}, \textbf{\code{Trades}}, \code{Fills}, \code{FillAggrs}, \textbf{\code{Universe}}, \textbf{\code{Violations}}, \code{ViolationsSummary}, \code{PerfCharts}\\\rowsep
\code{Contest} & \code{List}, \code{PfList}, \code{TopRankDetail}, \code{TopPick}\\\rowsep
\code{Analysis} & \code{PfAsOf}, \code{PrftByProd}, \code{HistPosChart}, \code{TrendChart}, \code{ActivityChart}\\\rowsep
\code{Util} & \code{SearchProducts}\\
\end{tabularx}}}{전체 경로는 영역과 이름을 이은 것이다(예: \code{Portfolio/Orders}). 굵은 글씨: 2026년 11월 24일 실제로 호출함(\code{PfBs}는 404, \figref{readeps}). 보통 글씨: 클라이언트 코드에서만 읽음. 주문 경로의 요청 본문은 \figref{orderapi}에 있다.}{웹 클라이언트 코드, DART Quant Team}


\lead{Bearer 헤더면 충분하며, 조회에는 쿠키가 필요 없다.} 실제 포트폴리오의 \code{Portfolio/Violations}에 \code{GET}을 보내면 헤더에 액세스 토큰이 있을 때 HTTP 200, 없을 때 401이 왔다.

\lead{포트폴리오 2개가 같은 대회 이름을 가질 수 있으므로, 운용 중인 포트폴리오는 대회 ID, 초기화 여부, 종료일로 골라야 한다.} \code{pfs} 목록은 대회에 참가하거나 초기화할 때마다 항목이 하나씩 늘어난다(\figref{pfs}). 13회차에는 대회 이름이 같은 항목이 2개 있어서, 이름으로 고르면 초기화로 끝난 10월 1\textasciitilde 8일 항목이 나올 수 있다.

% Korean edition of figs/pfs.tex.
\fig{pfs}{운용 중인 포트폴리오는 초기화되지 않았고 종료일도 지나지 않았다}[로그인 응답 \code{pfs}의 항목 6개, 최신순. 포트폴리오 ID는 생략]{%
{\tabfont\begin{tabularx}{\linewidth}{>{\RaggedRight\arraybackslash}p{34mm} >{\RaggedRight\arraybackslash}p{25mm} >{\RaggedRight\arraybackslash}p{13mm} >{\RaggedRight\arraybackslash}X}
\thd{대회 (\code{ctstId})} & \thd{기간} & \thd{\code{invalid}} & \thd{웹 페이지 상태}\\
\hl\textbf{13회차 (91)} & \textbf{10월 12일\textasciitilde 11월 30일} & \textbf{없음} & \textbf{운용 중}\\\rowsep
13회차 (91) & 10월 1일\textasciitilde 10월 8일 & True & 종료, 초기화\\\rowsep
12회차 (90) & 7월 3일\textasciitilde 8월 31일 & 없음 & 종료\\\rowsep
12회차 (90) & 7월 1일\textasciitilde 7월 2일 & True & 종료, 초기화\\\rowsep
11회차 (86) & 4월 1일\textasciitilde 5월 29일 & 없음 & 종료\\\rowsep
연습 (1) & 2월 5일\textasciitilde 4월 30일 & True & 종료, 초기화\\
\end{tabularx}}}{음영 행: 운용 중인 포트폴리오. \code{invalid}가 없고 기간이 11월 24일을 포함하는 유일한 항목이다. \code{invalid}는 초기화된 포트폴리오에만 나타난다. \code{active}는 운용 중인 것을 포함해 6개 모두 False이므로, 래퍼는 대회(\code{ctstId}), \code{invalid} 없음, 아직 지나지 않은 종료일(\code{lastD})로 고르며 정확히 1개가 맞아야 한다.}{DART Quant Team}


\lead{필드 이름 2개는 오해를 부른다.} \code{active}는 운용 중인 항목을 포함해 6개 모두 False이므로 어느 포트폴리오가 운용 중인지 알려주지 않는다. 초기화 여부는 \code{invalid}의 존재로 표시되므로, 키가 없으면 False로 읽어야 한다. 아래 규칙에 맞는 포트폴리오가 정확히 1개가 아니면 래퍼는 추측하지 않고 멈춘다.

\par\noindent\begin{minipage}{\linewidth}\fontsize{8.5}{11}\selectfont\ttfamily
\hspace*{4mm}select\_pf(ctstId):\\
\hspace*{8mm}cands $\leftarrow$ \{ pf $\in$ pfs : pf.ctstId = ctstId $\wedge$ $\neg$pf.invalid\\
\hspace*{38mm}$\wedge$ pf.lastD $\geq$ currBizDay \}\\
\hspace*{8mm}cands의 원소가 1개면 그것을 반환, 아니면 예외\end{minipage}\par\vspace{4pt}

\lead{날짜 매개변수는 어떤 날짜든 받고, 어느 날의 데이터를 돌려줄지는 서버가 정한다(\figref{dates}).} 주말 날짜는 직전 거래일, 미래 날짜는 최근 거래일 데이터를 돌려주고, 대회가 끝난 뒤의 휴일은 0을 돌려준다. 모두 HTTP 200이다. 따라서 날짜별로 이력을 모으는 호출자는 각 날짜를 거래일 달력과 대조해야 한다. 그러지 않으면 같은 날 데이터가 여러 날짜로 저장된다. 래퍼 자체에는 이런 반복문이 없다.

% Korean edition of figs/dates.tex.
\fig{dates}{주말과 미래 날짜는 오류 없이 다른 날의 데이터를 돌려준다}[요청 날짜별 \code{Portfolio/Violations} 응답, 2026년(서버 영업일: 11월 23일)]{%
{\tabfont\begin{tabularx}{\linewidth}{>{\RaggedRight\arraybackslash}p{44mm} R{12mm} R{11mm} >{\RaggedRight\setlength{\leftskip}{3mm}\arraybackslash}X}
\thd{요청 날짜} & \thd{HTTP} & \thd{바이트} & \thd{서버가 쓴 날짜와 근거}\\
11월 22일 (일요일) & 200 & 400 & 직전 거래일인 11월 20일: 소형주 비중이 3.21\%로, 11월 23일의 3.33\%와 다름\\\rowsep
11월 23일 (당일) & 200 & 400 & 11월 23일\\\rowsep
11월 24일 (다음 거래일) & 200 & 400 & 최근 거래일인 11월 23일: 11월 23일 응답과 바이트 단위로 같음\\\rowsep
12월 25일 (대회 이후 휴일) & 200 & 342 & 해당 날짜 없음: 모든 값이 0\\
\end{tabularx}}}{HTTP: 응답 상태 코드. 바이트: 응답 본문 크기. 응답은 \code{cmp}로 비교했고, 다른 경우 다른 필드를 출력했다. 요청 4개가 모두 성공하므로 호출자는 상태 코드로 데이터가 어느 날의 것인지 알 수 없다.}{DART Quant Team}


\lead{일요일 응답에는 11월 20일 데이터가 들어 있다. 11월 23일 응답과는 포트폴리오에 하나 있는 소형주의 비중만 다르다.} 그 종목의 가격은 11월 23일에 변하지 않았지만(전일 종가와 종가 모두 14,840원), 포트폴리오 가치는 그날 3.59\% 떨어졌다. 가격이 그대로인 종목은 포트폴리오가 줄어들면 비중이 커지므로, 11월 20일 비중은 3.33\,×\,(1\,−\,0.0359) = 3.21\%여야 하고, 일요일 요청이 돌려준 값이 바로 이것이다.

\lead{서버의 영업일은 자정 이후 몇 시간 동안 달력 날짜와 다르다.} 11월 24일 06:02:10에 \code{currBizDay}는 여전히 11월 23일이었고 07:06:44에는 11월 24일이었다. 즉 서버는 09:00 개장 전, 이 두 시각 사이에 다음 영업일로 넘어간다. 넘어가기 전에 보낸 요청에는 전날 날짜가 실린다.

\lead{조회 엔드포인트 3개가 더 있어 주문 가능 시간, 투자 유니버스, 주문 이력을 준다(\figref{readeps}).} \code{Session}은 주문이 체결되는 시간인 09:00\textasciitilde 15:20을 돌려준다. 15:20은 장 마감 동시호가가 시작되는 시각이다. \code{Universe}는 투자 가능 종목 2,648개와 GICS 섹터를 나열한다. 일부 종목코드는 한국거래소의 새 코드 체계에 따라 문자와 숫자가 섞여 있으므로(\code{A0015M0}) 종목코드는 문자열로 다룬다. \code{Orders}는 날짜 범위 \code{d0}, \code{d1}을 받는다. 처음에 \code{d}를 보냈을 때는 11월 23일에 대해 빈 목록이 왔지만, 실시간 피드에는 그날 체결된 주문이 있었다. 장중에 웹 페이지가 직접 보낸 요청은 \code{d0}과 \code{d1}을 보냈고, 완료된 주문을 포함해 범위 안의 모든 주문을 받았다(\ref{sec:orders}절).

% Korean edition of figs/readeps.tex.
\fig[wide]{readeps}{조회 엔드포인트 4개가 응답한다. \code{Orders}는 날짜 범위가 필요하고 \code{PfBs}는 404를 돌려준다}[운용 중인 \code{pfId}와 \code{d}로 호출한 \code{Portfolio} 조회 엔드포인트(07:33 KST), \code{d0}과 \code{d1}로 호출한 \code{Orders}(09:16 KST), 2026년 11월 24일]{%
{\tabfont\begin{tabularx}{\linewidth}{>{\RaggedRight\arraybackslash}p{14mm} R{10mm} R{16mm} >{\RaggedRight\setlength{\leftskip}{3mm}\arraybackslash}X >{\RaggedRight\arraybackslash}p{52mm}}
\thd{경로} & \thd{HTTP} & \thd{레코드} & \thd{내용} & \thd{래퍼에서의 용도}\\
\code{Session} & 200 & 1 & \code{hm0} 09:00, \code{hm1} 15:20 & 주문 시작·종료 시각의 범위\\\rowsep
\code{Universe} & 200 & 2,648 & \code{Id}, \code{nm}, \code{sec}(GICS 섹터). 일부 종목코드는 영숫자 혼합 & 종목코드 검증, 종목별 섹터\\\rowsep
\code{Trades} & 200 & 62 & 진입부터 청산까지 포지션당 레코드 1개. 07:33에 23개가 열려 있어 보유 종목 23개와 일치. 그날 진입 뒤 64개. 시각은 UTC & 포지션 이력과 실현 손익\\\rowsep
\code{Orders} & 200 & 주문당 1 & \code{d0}부터 \code{d1}까지의 모든 주문(\code{Done} 포함)이므로 개수는 보낸 날짜에 따라 다름. 필드: \code{Id}, \code{state}, \code{qty2ord}, \code{cumQty}, \code{avgPx}(평균 체결가, 원), \code{cnclT}, \code{genT}, \code{lastMCV}, \code{sgn}. \code{d}로 보내면 빈 목록 & 시작할 때 그날 주문 건수\\\rowsep
\code{PfBs} & 404 & 0 & 시도한 3개 영역(\code{Portfolio}, \code{Contest}, \code{Analysis}) 어디에도 없음, GET과 POST 모두 & 쓰지 않음. 잔고는 실시간 피드에서\\
\end{tabularx}}}{별도 표시가 없으면 쿼리 매개변수를 붙인 GET. HTTP: 상태 코드. 레코드: 응답의 항목 수. \code{Orders} 매개변수는 웹 페이지 자체의 요청에서 그대로 가져왔다.}{DART Quant Team}

\verdict{래퍼는 모든 요청의 날짜를 서버의 \code{currBizDay}로 채우며, \code{currBizDay}가 KST 달력 날짜와 다르면 주문을 거부한다(D6, D27). 포트폴리오는 위 규칙으로 고르고, 규칙에 맞는 포트폴리오가 없거나 둘 이상이면 멈춘다(D7).}
